\chapter{Urethrotomy}
\index{Urethrotomy}

\section*{Definition and Overview}
Urethrotomy is a specialized surgical procedure primarily indicated for the treatment of urethral strictures, which are abnormal narrowings of the urethral lumen caused by scarring of the lining epithelium and the surrounding corpus spongiosum (a condition known as spongiofibrosis). The urethra serves as the conduit for the excretion of urine from the urinary bladder to the external environment, and in males, it additionally serves as the pathway for semen during ejaculation. Structurally, any significant narrowing of this channel impedes the normal flow of urine, creating a mechanical obstruction that can lead to progressive lower urinary tract symptoms (LUTS), urinary retention, and in severe cases, upper urinary tract deterioration, renal impairment, or life-threatening urosepsis.

The term urethrotomy encapsulates several specific surgical techniques designed to incise the scar tissue obstructing the urethral lumen, thereby restoring its patency. In modern urological practice, the most common variant is Direct Vision Internal Urethrotomy (DVIU), also referred to as a Sachse urethrotomy. This procedure is performed endoscopically under direct visual guidance using a specialized instrument called a urethrotome. The surgeon inserts the sheath of the instrument into the urethra, visualizes the stricture, and uses a cold-knife blade (or a laser fiber) to make one or more longitudinal cuts through the scar tissue, allowing the urethra to expand. Another variation is the Otis urethrotomy, which is a blind internal urethrotomy traditionally used for female urethral calibration or to treat diffuse, mild narrowing of the male urethra by stretching and incising the tissue without direct endoscopic visualization. Historically, external urethrotomy was performed via an open perineal incision to divide strictures, though this has largely been superseded by endoscopic techniques and modern reconstructive urethroplasty.

Understanding the anatomical divisions of the urethra is crucial for the application of urethrotomy. In males, the urethra is divided into the anterior segment (comprising the meatus, fossa navicularis, penile or pendulous urethra, and bulbar urethra) and the posterior segment (comprising the membranous and prostatic urethra). Urethrotomy is primarily utilized for short, single strictures located in the bulbar urethra. Because the bulbar urethra is surrounded by a thick layer of corpus spongiosum, it can accommodate the longitudinal incision and the subsequent healing process, provided the scar tissue is not too extensive. Conversely, strictures in the penile urethra or those exceeding 1.5 centimeters in length are associated with high recurrence rates when treated with urethrotomy, often necessitating formal open reconstruction (urethroplasty). Urethrotomy remains a cornerstone of initial stricture management due to its minimally invasive nature, rapid recovery time, and immediate relief of obstructive symptoms, despite its notable risk of long-term stricture recurrence.

\section*{Epidemiology}
Urethral stricture disease, the primary pathology necessitating urethrotomy, represents a significant source of morbidity worldwide. The epidemiology of the condition is highly skewed toward the male population, owing to the greater length and anatomical complexity of the male urethra, making it far more susceptible to trauma, instrumentation, and infection compared to the short, relatively protected female urethra. The incidence of urethral strictures in males is estimated to be between 200 and 600 cases per 100,000 men, with the prevalence rising dramatically with age. While strictures can occur at any point in the lifespan, there is a clear bimodal distribution in presentation. A peak in young adulthood (ages 18 to 35) is typically associated with pelvic or perineal trauma, sports injuries, and sexually transmitted infections. A second, much larger peak occurs in older men (over the age of 65), driven primarily by iatrogenic factors, such as transurethral surgeries for benign prostatic hyperplasia (BPH) or prostate cancer, long-term indwelling urethral catheterization, and age-related ischemic changes.

In the female population, urethral strictures are exceedingly rare, accounting for less than 0.02\% of all urological complaints. The short length (approximately 4 centimeters) and high elasticity of the female urethra mean that narrowings are typically related to severe pelvic fractures, direct obstetric trauma during prolonged labor, or complications from anti-incontinence surgeries such as mid-urethral slings. Consequently, urethrotomy is rarely performed in females, and when it is, it is typically limited to gentle dilation or highly controlled calibration procedures like the Otis urethrotomy.

The socioeconomic burden of urethral strictures and their management through urethrotomy is substantial. In developed countries, hundreds of millions of dollars are expended annually on diagnostic imaging, endoscopic procedures, and post-operative monitoring. Urethrotomy remains the most frequently performed procedure for urethral stricture disease globally, largely due to its simplicity and low initial cost, even though the necessity for repeat interventions due to recurrence can accumulate significant long-term costs. In developing countries, the epidemiology of stricture disease differs, with a higher proportion of cases resulting from untreated sexually transmitted infections (such as gonorrhea) and traumatic pelvic injuries sustained in motor vehicle accidents or industrial incidents, where access to immediate reconstructive urology may be limited.

\section*{Aetiology and Pathophysiology}
The development of a urethral stricture requiring urethrotomy is a complex pathological process initiated by damage to the urethral epithelium and the underlying vascular structures. The causes and mechanisms can be categorized as follows:

\begin{itemize}
  \item \textbf{Iatrogenic Injury:} The most frequent cause of urethral stricture in modern clinical practice, accounting for approximately 30\% to 40\% of cases. Iatrogenic trauma occurs during medical interventions, including the insertion of indwelling urinary catheters, endoscopic instruments (such as cystoscopes, resectoscopes, and ureteroscopes), and surgeries like transurethral resection of the prostate (TURP) or holmium laser enucleation of the prostate (HoLEP). Prolonged catheterization can cause pressure necrosis of the urethral mucosa, particularly at the penoscrotal junction. The use of large-caliber instruments during transurethral surgery can cause frictional trauma, leading to mucosal tears and subsequent scarring. Additionally, radiation therapy for prostate or pelvic cancers can induce progressive endarteritis obliterans, resulting in tissue ischemia and dense, fibrotic strictures of the bulbar or membranous urethra.
  \item \textbf{Trauma:} Direct physical trauma is a major driver of urethral injury. Idiopathic or accidental trauma can take several forms. A classic ``straddle'' injury occurs when a person falls astride a hard object (such as a bicycle bar, fence, or ladder rung), compressing the bulbar urethra against the pubic symphysis, resulting in contusion or laceration. Severe pelvic fractures, often sustained in high-impact motor vehicle accidents, can shear the urethra at the membranous-bulbar junction, leading to pelvic fracture urethral distraction defects (PFUDD). Direct penetrating injuries (such as gunshot or knife wounds) and self-inflicted foreign body insertions are less common but highly damaging traumatic causes.
  \item \textbf{Inflammatory and Infectious Conditions:} Historically, infections were the leading cause of urethral strictures, and they remain a significant factor in regions with limited healthcare access. Urethritis caused by sexually transmitted pathogens, primarily \textit{Neisseria gonorrhoeae} and \textit{Chlamydia trachomatis}, leads to severe mucosal inflammation. If untreated or improperly treated, this inflammation can extend into the corpus spongiosum, causing chronic cicatrization and long, multi-focal strictures. Among non-infectious inflammatory conditions, Lichen Sclerosus (formerly known in males as balanitis xerotica obliterans or BXO) is a progressive, chronic inflammatory skin disease that affects the glans penis, prepuce, and distal urethra. It causes dense, white scarring that typically starts at the meatus and can migrate proximally to involve the entire anterior urethra (panurethral stricture).
  \item \textbf{Congenital Anomalies:} Congenital urethral strictures are rare and must be distinguished from posterior urethral valves. True congenital strictures are thought to arise from localized developmental failure of the urethral plate during embryogenesis, resulting in a thin, diaphragm-like narrowing, most commonly in the membranous or bulbar segments.
  \item \textbf{Idiopathic Causes:} In approximately 30\% of patients, no definitive cause can be identified. These are termed idiopathic strictures. It is hypothesized that many of these represent forgotten minor childhood trauma (such as a mild fall on a bicycle) or subclinical infections that slowly progressed over decades.
\end{itemize}

The pathophysiology of stricture formation centers on the concept of spongiofibrosis. The normal urethral mucosa is lined by pseudostratified columnar epithelium, which transitions to stratified squamous epithelium distally. This mucosa sits on a highly vascular, elastic supportive tissue layer within the corpus spongiosum. When the epithelial barrier is compromised by trauma, infection, or ischemia, urine is able to extravasate into the surrounding subepithelial tissues. The highly acidic and solute-rich nature of urine triggers a marked local inflammatory response. Fibroblasts migrate to the area, leading to the deposition of Type I collagen, which replaces the highly elastic Type III collagen and smooth muscle cells characteristic of normal corpus spongiosum. This fibrotic scar tissue contracts over time, reducing the elasticity of the urethral wall and narrowing the lumen. The compromised lumen creates a high-pressure zone proximal to the stricture during voiding. According to physical laws of fluid dynamics, the increased turbulence and shear stress at the margins of the stricture propagate further mechanical injury to the adjacent healthy epithelium, leading to proximal extension of the scar. This cyclic process explains why untreated strictures often worsen and elongate over time. Urethrotomy aims to break this cycle by dividing the ring of scar tissue, allowing the urethral lumen to expand to its physiological diameter.

\section*{Clinical Features}
The clinical presentation of urethral stricture disease is primarily characterized by progressive lower urinary tract symptoms (LUTS) resulting from the mechanical obstruction of urinary flow. The severity of symptoms correlates with the degree of luminal narrowing rather than the length of the stricture itself. The signs and symptoms can be categorized as follows:

\begin{itemize}
  \item \textbf{Obstructive Voiding Symptoms:} The classic presentation includes a progressive reduction in the caliber and force of the urinary stream. Patients often report that their urine stream has become thin, weak, or split (spraying). Hesitancy, where the patient must wait for the stream to initiate, and straining to void (Valsalva voiding) are common. Incomplete emptying is frequently reported, with patients feeling that urine remains in the bladder immediately after voiding. Post-void dribbling (the involuntary loss of urine immediately after finishing urination) occurs as trapped urine slowly drains from the dilated segment proximal to the stricture.
  \item \textbf{Irritative Voiding Symptoms:} Over time, the urinary bladder must generate higher intravesical pressures to overcome the urethral obstruction. This leads to detrusor muscle hypertrophy and instability, presenting as irritative symptoms. These include urinary frequency (needing to urinate more than eight times per day), urgency (a sudden, compelling desire to urinate that is difficult to defer), and nocturia (waking up multiple times during the night to void). Dysuria (painful urination) can occur, particularly if there is secondary mucosal inflammation or a co-existing urinary tract infection.
  \item \textbf{Hematuria and Discharge:} Microscopic or macroscopic hematuria (blood in the urine) may be present, resulting from the rupture of engorged capillaries in the inflamed mucosal lining or from the high-pressure passage of urine through the stricture site. A bloody or purulent urethral discharge may also be noted, especially in cases where the stricture is secondary to active urethritis or a periurethral abscess.
  \item \textbf{Ejaculatory Dysfunction:} In males, the urethra serves as the conduit for semen. A tight stricture can impede ejaculation, leading to a weak or dribbling ejaculatory stream, discomfort during ejaculation (odynophagia), or hematospermia (blood in the semen).
  \item \textbf{Recurrent Infections and Complications:} The stasis of urine in the bladder and the proximal urethra creates an environment highly conducive to bacterial proliferation. Patients may present with recurrent, symptomatic urinary tract infections (UTIs), acute epididymitis, or prostatitis. The retrograde transmission of infected urine through the ejaculatory ducts can cause severe scrotal pain, swelling, and fever.
  \item \textbf{Acute Urinary Retention:} In advanced cases, the urethral lumen may become completely occluded, leading to acute urinary retention (AUR). The patient experiences a sudden, painful inability to pass any urine, resulting in a visibly and palpably distended bladder in the suprapubic region. AUR is a medical emergency that requires immediate bladder decompression.
\end{itemize}

\section*{Differential Diagnosis}
When evaluating a patient with obstructive lower urinary tract symptoms, several alternative diagnoses must be considered, as many conditions can mimic the clinical presentation of a urethral stricture. The key differential diagnoses include:

\begin{itemize}
  \item \textbf{Benign Prostatic Hyperplasia (BPH):} This is the most common cause of bladder outlet obstruction in aging males. BPH involves the non-malignant hypertrophy of the transition zone of the prostate gland, which compresses the prostatic urethra. Unlike a urethral stricture, which can occur at any location along the urethra and affects men of all ages, BPH is localized to the prostatic segment and is rare in men under 50. Uroflowmetry and retrograde urethrography help differentiate the two.
  \item \textbf{Prostate Cancer:} Locally advanced adenocarcinoma of the prostate can invade the prostatic urethra or bladder neck, causing severe urinary obstruction. On clinical examination, a digital rectal exam (DRE) may reveal a hard, nodular, or asymmetric prostate, and serum Prostate-Specific Antigen (PSA) levels are frequently elevated.
  \item \textbf{Bladder Neck Contracture:} This is a localized scarring and narrowing of the bladder outlet, typically occurring as a complication of prior prostate surgeries, such as radical prostatectomy, TURP, or radiation therapy. While clinically similar to a bulbar urethral stricture, it is anatomically localized to the junction between the bladder and the prostatic urethra and is identified via cystourethroscopy.
  \item \textbf{Detrusor Sphincter Dyssynergia (DSD):} A neurological condition where the external urethral sphincter contracts involuntarily during detrusor contraction, rather than relaxing to allow urination. This is common in patients with spinal cord injuries or multiple sclerosis. Urodynamic studies are essential to distinguish this functional obstruction from a mechanical stricture.
  \item \textbf{Urethral Carcinoma:} Primary malignancy of the urethra is extremely rare but must be excluded, particularly in patients with rapid-onset obstruction, perineal mass, hematuria, or persistent urethral discharge. Biopsy during cystoscopy provides a definitive diagnosis.
  \item \textbf{Urethral Foreign Body:} The insertion of foreign objects into the urethra, either accidentally or intentionally, can cause sudden obstruction and severe pain. Diagnosis is confirmed through clinical history, plain radiography, or direct visualization.
  \item \textbf{Meatal Stenosis:} A narrowing of the external urethral meatus, often caused by chronic inflammation, Lichen Sclerosus, or complications of circumcision. It is easily diagnosed on physical examination by inspecting the tip of the penis.
\end{itemize}

\section*{When to Seek Medical Care}
Patients experiencing changes in their urinary patterns should seek prompt medical evaluation. While a gradual reduction in urinary stream is a common symptom that should be discussed with a general practitioner or urologist during a scheduled appointment, certain acute signs and symptoms require emergency medical care.

Immediate emergency medical attention should be sought if the patient experiences any of the following:
\begin{itemize}
  \item \textbf{Complete Inability to Urinate (Anuria):} A sudden, complete cessation of urine flow accompanied by intense, agonizing pain in the lower abdomen. This indicates acute urinary retention, which can lead to bladder rupture or acute kidney injury if not decompressed promptly.
  \item \textbf{Signs of Systemic Infection (Urosepsis):} High fever, chills, rapid heart rate (tachycardia), rapid breathing (tachypnea), confusion, or extreme lethargy. Urosepsis is a life-threatening condition resulting from infected urine entering the bloodstream under high pressure behind an obstruction.
  \item \textbf{Severe Hematuria:} The passage of large amounts of bright red blood or thick blood clots in the urine, which can lead to urinary tract obstruction (clot retention) and acute anemia.
  \item \textbf{Severe Scrotal or Perineal Pain and Swelling:} Sudden, intense pain accompanied by redness and swelling in the scrotum or the area between the scrotum and anus, which may indicate acute epididymitis or a periurethral abscess/extravasation.
\end{itemize}

For life-threatening symptoms, patients should immediately contact their local emergency services (for example, \textbf{local emergency services} in many countries, \textbf{911} in the United States and Canada, or \textbf{999} in the United Kingdom) or proceed to the nearest hospital emergency department.

\section*{Diagnosis and Investigations}
Confirming the diagnosis of a urethral stricture and planning for a urethrotomy requires a systematic combination of clinical evaluation, functional flow studies, and anatomical imaging. The standard diagnostic workup includes:

\begin{itemize}
  \item \textbf{Clinical History and Physical Examination:} The initial step involves a detailed review of symptoms, prior urological procedures, catheterizations, trauma history, and sexual history. Physical examination includes palpation of the ventral aspect of the penis and scrotum to detect areas of induration or scarring along the course of the urethra. The external urethral meatus is inspected to rule out meatal stenosis or active inflammatory skin disease like Lichen Sclerosus.
  \item \textbf{Uroflowmetry:} A non-invasive screening test where the patient voids into a specialized funnel connected to a flow meter. The device records the volume of urine passed per second, generating a flow curve. A normal flow curve is bell-shaped, with a peak flow rate exceeding 15 milliliters per second (mL/s). In the presence of a urethral stricture, the flow curve displays a characteristic ``plateau'' pattern, where the flow rate remains low (peak flow rate often less than 10 mL/s) throughout an elongated voiding cycle.
  \item \textbf{Post-Void Residual (PVR) Measurement:} Performed immediately after uroflowmetry using a transabdominal ultrasound scanner. A PVR volume greater than 50 to 100 mL indicates significant bladder emptying failure, highlighting the severity of the obstruction and the potential risk of upper urinary tract distress.
  \item \textbf{Retrograde Urethrogram (RUG):} The gold standard for anatomical imaging of the anterior urethra. A radiopaque contrast medium is slowly injected into the distal urethra via a small catheter tip inserted into the meatus. Fluoroscopic X-ray images are taken as the contrast fills the urethral channel. RUG accurately defines the location, number, length, and caliber of strictures in the penile and bulbar segments.
  \item \textbf{Voiding Cystourethrogram (VCUG):} Often performed in conjunction with a RUG to evaluate the posterior urethra. The bladder is filled with contrast medium, and fluoroscopic images are captured while the patient voids. This is essential for visualizing the bladder neck and prostatic urethra, which cannot be fully evaluated with a retrograde study alone.
  \item \textbf{Cystourethroscopy (Diagnostic):} Endoscopic examination of the urethra using a flexible or rigid cystoscope under local anesthesia. This allows direct visualization of the stricture, assessment of the mucosal quality, and evaluation of the caliber of the narrowing. In cases where the stricture is too tight to allow the passage of the scope, a pediatric cystoscope or a guide wire may be used to assess the distal margin.
  \item \textbf{Urinalysis and Urine Culture:} Mandatory prior to any diagnostic instrumentation or surgical intervention. An active, untreated urinary tract infection is an absolute contraindication to urethroscopy and urethrotomy, as instrumenting an infected system can precipitate urosepsis. If bacteria are identified, a targeted course of antibiotics must be completed before the procedure.
\end{itemize}

\section*{Management}
The management of urethral stricture disease through urethrotomy involves precise surgical execution and structured post-operative care, balancing immediate symptom relief with the long-term risk of recurrence.

\begin{itemize}
  \item \textbf{Surgical Options for Urethrotomy:}
  \begin{itemize}
    \item \textit{Direct Vision Internal Urethrotomy (DVIU):} The patient is typically placed in the lithotomy position under general, spinal, or local anesthesia. The surgeon inserts a specialized urethrotome equipped with a cold steel knife blade. Under direct visual guidance, the knife is advanced, and one or more longitudinal incisions are made through the stricture ring. The incision is typically made at the 12 o'clock position (dorsal aspect) to avoid damaging the highly vascular corpus spongiosum on the ventral and lateral sides, which could lead to significant hemorrhage. The cut is made deep enough to expose healthy, vascular tissue beneath the scar, allowing the urethral lumen to expand to approximately 20 to 24 French caliber.
    \item \textit{Laser Urethrotomy:} Similar to DVIU, but the mechanical knife is replaced with a laser fiber, most commonly using a Holmium:YAG or Thulium laser. The laser energy is used to vaporize or cut the scar tissue. While laser urethrotomy offers excellent intraoperative hemostasis, clinical trials have shown that long-term recurrence rates are comparable to cold-knife DVIU.
    \item \textit{Otis Urethrotomy:} A procedure where a graduated Otis urethrotome is inserted into the urethra, dilated to a predetermined caliber, and a knife blade is drawn along the instrument to make a longitudinal incision. This technique is reserved for female urethral strictures, diffuse calibration, or as a preliminary step to facilitate the passage of large-caliber instruments (like resectoscopes) during transurethral surgeries.
  \end{itemize}
  \item \textbf{Post-operative Care and Catheterization:} Following the procedure, a Foley catheter is placed through the freshly cut urethra into the bladder. The catheter acts as a temporary stent, keeping the incised edges of the urethra apart and allowing epithelialization to occur over the raw surface, theoretically preventing the wound edges from fusing back together. The optimal duration of catheterization remains a subject of debate, ranging from 3 days to 2 weeks, depending on the severity of the stricture and the surgeon's preference. Longer duration does not consistently correlate with lower recurrence rates but may increase the risk of catheter-associated UTIs.
  \item \textbf{Adjuvant Therapies and Self-Dilation:} To reduce the risk of stricture recurrence, patients are often instructed on Clean Intermittent Self-Catheterization (CISC) or Clean Intermittent Self-Dilatation (CISD) after catheter removal. The patient is taught to gently insert a lubricated, single-use catheter into the urethra once or twice a week for several months. This serves to mechanically disrupt any early scar formation. Additionally, the instillation of topical pharmacological agents, such as steroids (triamcinolone) or chemotherapeutic agents (mitomycin C), into the incision site or via the catheter has been investigated to inhibit fibroblast proliferation, showing some promise in reducing recurrence rates.
  \item \textbf{Surgical Comparison and Limitations:} Urethrotomy is highly successful for short, single, bulbar strictures (less than 1.5 cm in length) without extensive spongiofibrosis, yielding success rates of 50\% to 70\% at one year. However, for strictures exceeding 2 cm, penile urethral strictures, or cases where urethrotomy has already failed once, the success rate of a repeat urethrotomy drops to less than 10\% to 20\%. In these scenarios, reconstructive surgery (urethroplasty), which involves excising the scar tissue and primary anastomosis or substituting the tissue with a graft (such as buccal mucosa), is the treatment of choice, offering long-term success rates exceeding 85\% to 90\%.
\end{itemize}

\section*{Complications}
Although urethrotomy is considered a minimally invasive and relatively safe procedure, it carries inherent risks of short-term and long-term complications. These include:

\begin{itemize}
  \item \textbf{Intraoperative Complications:}
  \begin{itemize}
    \item \textit{Bleeding/Hemorrhage:} The corpus spongiosum surrounding the urethra is a highly vascular erectile tissue. Incising too deeply or off-center can result in significant bleeding. While minor bleeding is common and typically managed with catheter balloon traction or compression, severe hemorrhage may require packing or formal surgical exploration.
    \item \textit{Fluid Extravasation:} During the procedure, irrigation fluid (typically sterile saline) is used to maintain visibility. If the incision extends through the corpus spongiosum into the deep perineal spaces, irrigation fluid can extravasate into the scrotum, perineum, or lower abdominal wall, leading to severe edema and potential tissue necrosis.
    \item \textit{False Passage:} The blind or difficult manipulation of instruments in a narrow urethra can result in the creation of a ``false passage'' where the instrument dissects a channel parallel to the true urethral lumen.
  \end{itemize}
  \item \textbf{Early Postoperative Complications:}
  \begin{itemize}
    \item \textit{Urinary Tract Infection and Epididymitis:} Bacterial colonization of the catheter or the surgical site can lead to cystitis, pyelonephritis, or acute epididymo-orchitis, requiring systemic antibiotic therapy.
    \item \textit{Urinary Retention:} If the post-operative catheter becomes blocked with blood clots or is expelled prematurely, the patient may experience painful urinary retention.
  \end{itemize}
  \item \textbf{Late Postoperative Complications:}
  \begin{itemize}
    \item \textit{Stricture Recurrence:} The most common long-term complication. The incision created during urethrotomy heals by secondary intention, which can lead to recurrent fibrotic scarring and stricture formation, often worse or longer than the original stricture.
    \item \textit{Erectile Dysfunction (ED):} Rare, but can occur if the dorsal incision at the 12 o'clock position damages the cavernous nerves or if deep lateral cuts compromise the vascular supply to the corpora cavernosa.
    \item \textit{Urinary Incontinence:} If the stricture is located near the membranous urethra, the incision may injure the external urethral sphincter, leading to stress urinary incontinence.
    \item \textit{Urethral Diverticulum or Fistula:} Deep incisions and high pressure voiding can lead to outpouching of the urethral wall (diverticulum) or the formation of an abnormal communication to the skin (fistula).
  \end{itemize}
\end{itemize}

\section*{Prognosis and Outcomes}
The prognosis for patients undergoing urethrotomy depends heavily on stricture characteristics, including location, length, etiology, and the presence of prior interventions. For a patient presenting with their first, single, short (less than 1.5 cm) stricture located in the bulbar urethra, the immediate success rate of DVIU is favorable, with approximately 50\% to 70\% of patients remaining free of recurrence at one year, and about 30\% to 40\% at five years.

However, the prognosis declines precipitously with repeat urethrotomies. Clinical studies have conclusively shown that if a primary urethrotomy fails, a second urethrotomy has a long-term success rate of less than 20\%. Third and subsequent urethrotomies have a success rate approaching 0\%, as each sequential incision adds to the cumulative cicatrization and spongiofibrosis, making the stricture denser and longer. Therefore, guidelines from international urological societies (such as the American Urological Association and the European Association of Urology) recommend that a patient should not undergo more than one or two internal urethrotomies before being referred for definitive reconstructive urethroplasty.

The recovery time following a urethrotomy is rapid. Most patients are discharged on the day of surgery or the following morning. Mild hematuria and perineal discomfort are common for the first few days and are managed with oral analgesics. Once the urinary catheter is removed (typically 3 to 10 days post-operatively), the patient experiences immediate improvement in the force of the urinary stream. Follow-up monitoring is critical and usually involves uroflowmetry and symptom score questionnaires at 3, 6, and 12 months, and annually thereafter, to detect early signs of recurrence. If a patient remains recurrence-free for five years, the likelihood of long-term cure is high, though late recurrences can still occur.

\section*{Prevention and Self-Care}
Preventing the initial development of a urethral stricture and reducing the likelihood of recurrence after a urethrotomy involves a combination of medical vigilance and proactive self-care.

Key prevention and self-care measures include:
\begin{itemize}
  \item \textbf{Prevention of Initial Urethral Injury:}
  \begin{itemize}
    \item \textit{Careful Urinary Catheterization:} Healthcare providers must employ aseptic, gentle techniques and utilize generous amounts of sterile lubricant when inserting urinary catheters. Using the smallest appropriate catheter size helps minimize mucosal friction.
    \item \textit{Prompt Treatment of Infections:} Sexually active individuals should practice safe sex to prevent urethritis. Any symptoms of urethral discharge, burning during urination, or genital sores should be promptly evaluated and treated with appropriate antibiotics to prevent chronic, scarring inflammation.
    \item \textit{Safety and Protective Gear:} Wearing appropriate protective athletic gear (such as cups or padded shorts) during high-impact or contact sports can help prevent traumatic straddle injuries to the perineum.
  \end{itemize}
  \item \textbf{Post-operative Self-Care after Urethrotomy:}
  \begin{itemize}
    \item \textit{Catheter Management:} Patients must maintain strict hygiene while the urinary catheter is in place. The catheter bag should be kept below the level of the bladder to prevent retrograde flow. The insertion site at the urethral meatus should be cleaned daily with mild soap and water.
    \item \textit{Adherence to Self-Dilation:} If prescribed, performing Clean Intermittent Self-Dilatation (CISD) precisely as instructed by the urologist. Consistency in inserting the dilation catheter helps keep the healing incision open and prevents premature fusion of the wound edges.
    \item \textit{Adequate Hydration:} Drinking plenty of water (at least 2 liters per day, unless contraindicated by other medical conditions) helps dilute the urine, reducing irritation to the healing urethral tissues and minimizing the risk of urinary tract infections.
    \item \textit{Avoiding Strenuous Activity:} Refraining from heavy lifting, vigorous exercise, sexual intercourse, and activities that place direct pressure on the perineum (such as cycling or horseback riding) for at least 2 to 4 weeks after the procedure to prevent bleeding and disruption of the surgical site.
  \end{itemize}
\end{itemize}

\section*{Sources and Further Reading}
For reliable, evidence-based medical information on urethrotomy and urethral stricture disease, the following resources are recommended:

\begin{itemize}
  \item \textbf{American Urological Association (AUA):} The AUA provides comprehensive clinical guidelines for the diagnosis and treatment of urethral strictures. Urologists worldwide refer to these guidelines for evidence-based management pathways. Detailed patient education resources are available on their official portal: \texttt{https://www.urologyhealth.org}.
  \item \textbf{European Association of Urology (EAU):} The EAU publishes annually updated guidelines on urethral strictures, offering detailed insights into surgical options, outcomes, and follow-up protocols. The patient-facing guides can be accessed via: \texttt{https://patients.uroweb.org}.
  \item \textbf{National Institute of Diabetes and Digestive and Kidney Diseases (NIDDK):} Part of the U.S. National Institutes of Health (NIH), the NIDDK provides authoritative, easy-to-understand explanations of urinary tract obstruction, diagnostic tests, and treatment options: \texttt{https://www.niddk.nih.gov}.
  \item \textbf{Urology Care Foundation:} The official foundation of the AUA, offering peer-reviewed brochures, videos, and articles detailing what to expect before, during, and after urethroscopy and urethrotomy: \texttt{https://www.urologyhealth.org/urology-a-z/u/urethral-stricture-disease}.
  \item \textbf{Healthy WA (Government of Western Australia Department of Health):} Provides highly practical clinical guidelines and post-operative self-care instructions for patients recovering from internal urethrotomy and managing urinary catheters: \texttt{https://www.healthywa.wa.gov.au/Articles/U-Z/Urethral-strictures}.
  \item \textbf{NHS Choices (National Health Service, UK):} Offers comprehensive overviews of urethral stricture treatments, patient pathways, risks, and recovery guidelines within the public health system: \texttt{https://www.nhs.uk/conditions/urethral-stricture}.
\end{itemize}